In earlier sections, when discussing automated invariant placement and establishment of semantically valid narratives, I was discussing a form of semantic reasoning over the closed semantic world. While this is particularly relevant to automated failure diagnosis (through elimination of known semantically invalid narratives), ultimately, any form of reasoning over the closed semantic world can be achieved given that the necessary meaning has been established. 

Through this process, the apparent complexity of software systems will be collapsed into the clarity of a coherent design. Since the design is a closed semantic world, the ability to ask any meaningful question about that world and to be able to automatically receive meaningful answers to automate the management of that world is a very powerful capability.

For example, the engine can answer if the requirements which motivated a design are satisfied by the closed semantic world. This is possible because the requirements have shared meaning with the design and this enables the engine to automatically answer if it is satisfied. Another example could be the elimination of redundant steps in the narrative by automatically identifying if particular paths or subpaths have no meaningful effect on the world. These are just a few examples but the type of questions that can be asked is simply limited by the world meaning that has been established through the CSM.

This does not invent anything new, it is simply formalizing the existing software engineering process into an unambiguous and computable model so that it can be automated. The novelty of the CSM is that it has made the automation of software systems possible by enabling the engine to reason about the meaning of the closed semantic world in the same way that humans do. This is made possible by constructing the world from computable semantic primitives that eliminate the semantic gap between the implementation and the design. In this sense, this practical mechanism to automate the management of software systems is simply the logical conclusion of what was already common knowledge. 

When the CSM is used to automate the management of software systems, it is evident that through semantic reasoning and automated failure diagnosis a deterministic learning loop emerges to expand the design semantics through root cause analysis on environments that caused the failure. This is a meaningful part of fully automating systems management because the meaning of the failure is unambiguous and it results in a deterministic process that motivates the management steps. However, to enable this deterministic process at scale, every environment that caused a failure must be preserved.

To solve this problem, the unambiguous domain structure of the design is exploited in multiple ways. First, the minimal information necessary to replay the execution is unambiguously established by specifying all the information that the design cannot deterministically reproduce. Second, the unambiguous domain structure of the data is exploited to apply domain specific compression to minimize the cost of preserving the environments. To make this possible, an open source tool named Compressed Log Processor is leveraged. It has proven that domain specific compression is practical at a petabyte scale.

The novelty of CLP is that as part of the automated management, the environments that caused failures must be preserved at scale to enable the design to learn new semantics through root cause analysis. Without this, the learning loop would be broken and the management would not be automatic. Thus, by exploiting the unambiguous domain specification of the CSM, CLP enables the automation of software systems management by preserving the environments that teach the design new semantics through root cause analysis at scale.

In many ways, from a research perspective, CLP is the real novel solution that made automated software systems management possible. While the CSM does establish functionality that didn't exist before, the actual mechanism that would make the automated reasoning possible was already clear. It would involve formalizing the existing software engineering process into an unambiguous and computable model so that it can be automated. However, the same could not be said about CLP and in establishing domain specific compression as a proven technique, it has established a previously unknown capability that is necessary to make automated software systems management practical. 